\DocumentMetadata{
  lang = pt-BR,
  pdfstandard = A-2b,
  pdfversion = 1.7
}

\documentclass{abntexto-ufc}

\makeatletter
\newcommand{\ufcLongQuoteProbe}{%
  \typeout{UFC-LONGQUOTE-HANG=\the\hangindent}%
  \typeout{UFC-LONGQUOTE-FONTSIZE=\f@size}%
  \typeout{UFC-LONGQUOTE-BASELINE=\the\baselineskip}%
}
\makeatother

\begin{document}
\textual

\section{Estruturas normativas complementares}

Texto anterior à citação direta longa.

\Enquote{\ufcLongQuoteProbe
Citação direta longa de validação normativa, deliberadamente extensa para ocupar mais de uma linha e exercitar o tamanho reduzido, o espaçamento simples e o recuo de quatro centímetros fornecidos pela classe-base.}

Texto posterior à citação direta longa.

\begin{ufclettereditems}
  \item Alínea normativa de primeiro nível;
  \item Segunda alínea normativa:
    \begin{ufcdashedsubitems}
      \item Subalínea normativa de segundo nível.
    \end{ufcdashedsubitems}
\end{ufclettereditems}

A Equação~\ref{eq:normativa} exercita a numeração de fórmulas.
\begin{equation}
  E = mc^2
  \label{eq:normativa}
\end{equation}

\appendix{Instrumento de validação normativa}
Conteúdo do apêndice complementar.

\annex{Documento externo de validação normativa}
Conteúdo do anexo complementar.\footnote{Referência específica do anexo apresentada no próprio anexo para exercitar o requisito condicional.}

\end{document}
